The Linear Representation Hypothesis (\lrh) posits that high-level concepts are encoded as linear directions in the residual stream of Large Language Models (LLMs). While foundational work has validated this hypothesis for isolated concepts, its robustness under composition---the combination of multiple concept directions---remains poorly understood. In this work, we systematically investigate the compositionality of linear directions in \qwen. Contrary to our initial hypothesis that related concepts would compose more naturally, we uncover a \textbf{\consistencyparadox}: directions for \textit{unrelated} or novel concept combinations are significantly more consistent across contexts (Avg. Cosine Similarity $\sim$0.65) than those for common \textit{related} concepts (Avg. Cosine Similarity $\sim$0.58). Furthermore, we demonstrate that steering interventions using unrelated directions generalize as effectively as, and in some cases outperform, steering with related directions. Our results suggest that while LLMs utilize linear superposition as a default mechanism for novel compositions, they develop specialized non-linear representations for frequently co-occurring concepts. This finding challenges the assumption that linearity is the universal ideal for concept representation and suggests that the failure of linear compositionality in related concepts is a hallmark of feature "compilation" into more efficient non-linear manifolds.
